\documentclass{article}

\usepackage{amsmath,amssymb}
\usepackage{xurl}
\usepackage[hidelinks]{hyperref}
\setlength{\emergencystretch}{2em}

% Shared commands for notes in docs/paper-gaps/.

\DeclareUnicodeCharacter{2080}{\ifmmode{}_{0}\else\textsubscript{0}\fi}
\DeclareUnicodeCharacter{2081}{\ifmmode{}_{1}\else\textsubscript{1}\fi}
\DeclareUnicodeCharacter{2082}{\ifmmode{}_{2}\else\textsubscript{2}\fi}
\DeclareUnicodeCharacter{2099}{\ifmmode{}_{n}\else\textsubscript{n}\fi}

\newcommand{\Lean}{\textsc{Lean}}
\ExplSyntaxOn
\NewDocumentCommand{\leanid}{m}{
  \ifmmode
    \textsf{\detokenize{#1}}
  \else
    \group_begin:
    \urlstyle{sf}
    \tl_set:Nn \l_tmpa_tl {#1}
    \tl_replace_all:Nnn \l_tmpa_tl {\_} {_}
    \exp_args:NV \path \l_tmpa_tl
    \group_end:
  \fi
}
\ExplSyntaxOff
\providecommand{\tr}{\operatorname{tr}}
\newcommand{\tnleanIssueBase}{https://github.com/LionSR/TNLean/issues/}
\newcommand{\tnleanPrBase}{https://github.com/LionSR/TNLean/pull/}
\newcommand{\tnleanDocBase}{https://sirui-lu.com/TNLean/docs/}
\newcommand{\ghissue}[1]{\href{\tnleanIssueBase#1}{GitHub issue \##1}}
\newcommand{\ghpr}[1]{\href{\tnleanPrBase#1}{PR \##1}}
\newcommand{\doclink}[1]{\href{\tnleanDocBase#1.html}{\path{#1}}}

% Dirac notation and the identity operator, as in blueprint/src/macros/common.tex.
% \providecommand keeps note-local definitions authoritative.
\usepackage{dsfont}
\providecommand{\ket}[1]{|#1\rangle}
\providecommand{\bra}[1]{\langle #1|}
\providecommand{\braket}[2]{\langle #1 | #2 \rangle}
\providecommand{\ketbra}[2]{|#1\rangle\!\langle #2|}
\providecommand{\Id}{\mathds{1}}
\providecommand{\one}{\mathds{1}}
\providecommand{\C}{\mathbb{C}}
\providecommand{\MN}[1]{M_{#1}(\C)}
\providecommand{\MD}{M_D(\C)}

% External referencing for paper sources.  The build helper writes sanitized
% auxiliary files under build/paper-aux/ and excludes duplicated source labels.
\usepackage{xr}

% Import a generated paper auxiliary file with a caller-supplied label prefix.
% The candidate paths support builds from docs/paper-gaps/, the repository root,
% and build/paper-aux/ itself.
\newcommand{\importpaperaux}[2]{%
  \IfFileExists{../../build/paper-aux/#2.aux}{%
    \externaldocument[#1]{../../build/paper-aux/#2}%
  }{%
    \IfFileExists{build/paper-aux/#2.aux}{%
      \externaldocument[#1]{build/paper-aux/#2}%
    }{%
      \IfFileExists{#2.aux}{%
        \externaldocument[#1]{#2}%
      }{}%
    }%
  }%
}

% General external reference macros
\newcommand{\paperref}[2]{\ref{#1-#2}}
\newcommand{\papereqref}[2]{(\ref{#1-#2})}

% CPSV16 source (1606.00608 / MPDO-22-12-17-2.tex) external document and shortcuts
\importpaperaux{cpsv16-}{1606.00608/MPDO-22-12-17-2-tnlean}
\newcommand{\cpsvref}[1]{\paperref{cpsv16}{#1}}
\newcommand{\cpsveqref}[1]{\papereqref{cpsv16}{#1}}

% Verdict marker: \gapnote{<kind>}{<status>}. Kind is the policy
% classification (clarification, local-correction, scope-restriction,
% unfaithful, false-source, open-gap); status is open, wip (elimination
% underway), resolved, or historical. Machine-read by the site index;
% prints a verdict line.
\newcommand{\gapnote}[2]{\par\noindent\textbf{Verdict:} #1 (#2).\par}


\title{Two-Projection Decomposition: Repeated-Projector Typo}
\author{}
\date{2026-08-12}

\begin{document}
\maketitle
\gapnote{local-correction}{resolved}

\section{Source passage}

The proof of Proposition~IV.5 in Ref.~\cite{Cirac2017MPU} applies the
two-projection lemma to projectors $P$ and $Q$ on $\mathbb{C}^D$. At
lines~760--763 of the local source file, the printed decomposition is
\begin{align}
    P
        &= \bigoplus_i \ketbra{0}{0}_i\oplus R,
    \label{eq:mpu_two_projection_projector_typo_printed_first_projector}\\
    P
        &= \bigoplus_i \ketbra{v_i}{v_i}_i\oplus S.
    \label{eq:mpu_two_projection_projector_typo_printed_repeated_projector}
\end{align}
The second occurrence of $P$ in
Eq.~\ref{eq:mpu_two_projection_projector_typo_printed_repeated_projector} is a
typographical error. The corrected decomposition is
\begin{align}
    P
        &= \bigoplus_i \ketbra{0}{0}_i\oplus R,
    \label{eq:mpu_two_projection_projector_typo_corrected_first_projector}\\
    Q
        &= \bigoplus_i \ketbra{v_i}{v_i}_i\oplus S.
    \label{eq:mpu_two_projection_projector_typo_corrected_second_projector}
\end{align}
The preceding sentence in Ref.~\cite{Cirac2017MPU} introduces the two
projectors as $P$ and $Q$, and the subsequent calculation uses $PQ$ and $QP$.
These surrounding statements determine the correction in
Eq.~\ref{eq:mpu_two_projection_projector_typo_corrected_second_projector}.

\section{Formalized interpretation}

The declarations
\leanid{LinearMap.IsSymmetricProjection.angleDefect_block} and
\leanid{LinearMap.IsSymmetricProjection.exists_orthonormal_angle_block} use the
two symmetric projections $P$ and $Q$. They establish the exact action of both
operators on a generic two-dimensional angle block.

The declarations
\leanid{LinearMap.IsSymmetricProjection.angleDefectSum} and
\leanid{LinearMap.IsSymmetricProjection.angleDefectComplement} assemble the
compression-indexed angle blocks into an orthogonal internal sum and its
invariant orthogonal complement. The restricted projections commute on that
complement. This construction includes the endpoint coordinates and implements
the corrected two-projector interpretation without adding a hypothesis or
altering the mathematical claim.

The declaration
\leanid{LinearMap.IsSymmetricProjection.reducedProjection} defines
$\widetilde P$ as the orthogonal projector onto $\operatorname{ran}(PQ)$. The
source defines this projector blockwise. Identifying its blockwise projector
with the canonical coordinate-free construction remains part of
\ghissue{6294}. No such identification is required for the formalized
absorption and range identities.

The associated theorems prove the three absorption identities from
lines~767--769 of Ref.~\cite{Cirac2017MPU}. They also prove
\begin{align}
    \operatorname{ran}(\widetilde P)
        &= \operatorname{ran}(PQ),
    \label{eq:mpu_two_projection_projector_typo_first_range_identity}\\
    \operatorname{ran}(PQP)
        &= \operatorname{ran}(PQ),
    \label{eq:mpu_two_projection_projector_typo_second_range_identity}\\
    \operatorname{ran}(\widetilde P Q)
        &= \operatorname{ran}(\widetilde P).
    \label{eq:mpu_two_projection_projector_typo_third_range_identity}
\end{align}
The final identity,
Eq.~\ref{eq:mpu_two_projection_projector_typo_third_range_identity}, supplies
the derived hypothesis for the equal-range right-factor theorem, which gives
the invertible map used at line~770 of Ref.~\cite{Cirac2017MPU}.

\section{Verdict}

\textbf{Verdict: resolved local fix.}
The repeated $P$ in the second decomposition is a resolved local
typographical error. Replacing it by $Q$ gives the intended angle-block
decomposition and leaves the mathematical statement unchanged.

The formal development uses $Q$ as the second projector, assembles all
compression-indexed blocks into an orthogonal internal sum, and supplies a
commuting invariant complement. It also formalizes the reduced projector and
the absorption and range identities used at lines~764--770 of
Ref.~\cite{Cirac2017MPU}. The canonical coordinate-free reduced projector has
not been identified with the source's particular blockwise projector, but that
identification is unnecessary for the proved identities.

\bibliographystyle{alpha}
\bibliography{references}

\end{document}